\chapter{Producto y alcance}

\section{Propósito}

SassBlum es una plataforma web para registrar, asignar, atender y auditar solicitudes de servicio
técnico. Sustituye el seguimiento disperso por un flujo único en el que cada ticket conserva
cliente, responsable, prioridad, estado, comentarios, adjuntos e historial.

El sistema combina un sitio público, autenticación y paneles separados para Cliente, Trabajador y
Administrador. También ofrece notificaciones, catálogo de servicios, contenido institucional,
testimonios moderados y reportes exportables en PDF y Excel.

\section{Valor entregado}

\begin{tabularx}{\textwidth}{@{}L{0.25\textwidth}Y Y@{}}
\toprule
\textbf{Necesidad} & \textbf{Respuesta del sistema} & \textbf{Resultado operativo}\\
\midrule
Ingreso de solicitudes & Formulario de ticket con servicio, prioridad, descripción y adjuntos & Entrada uniforme y consultable\\
Responsabilidad & Asignación y reasignación por administradores & Responsable visible y trazable\\
Seguimiento & Estados, comentarios e historial de eventos & Menos pérdida de contexto\\
Información oportuna & Notificaciones en aplicación, correo y WebSocket según configuración & Actualizaciones por el canal disponible\\
Gestión & Indicadores y exportación PDF/Excel & Seguimiento de volumen y estado\\
\bottomrule
\end{tabularx}

\section{Roles}

\begin{description}
  \item[Cliente.] Se registra, verifica su correo, crea tickets y consulta únicamente los propios.
  Puede revisar historial, notificaciones y administrar su testimonio.
  \item[Trabajador.] Consulta los tickets que le corresponden, añade comentarios y cambia estados
  operativos con una explicación obligatoria.
  \item[Administrador.] Gestiona usuarios, contenido, asignaciones, todos los tickets y reportes.
\end{description}

La interfaz oculta acciones no permitidas, pero la autorización definitiva siempre ocurre en el
backend. El aislamiento de tickets se implementa mediante permisos y consultas filtradas; no debe
describirse como RLS general de Supabase.

\section{Ciclo de vida del ticket}

\begin{center}
\fbox{\parbox{0.88\textwidth}{
\centering\sffamily
Nuevo $\xrightarrow{\text{asignación + comentario}}$ EnProceso
$\leftrightarrow$ EnEspera $\leftrightarrow$ Resuelto $\leftrightarrow$ Cerrado
}}
\end{center}

Un ticket \textbf{Nuevo} solo entra en trabajo cuando un administrador lo asigna. El personal
autorizado puede moverlo entre los cuatro estados operativos. \textbf{Cerrado no es terminal}: puede
reabrirse conservando la trazabilidad. Cada transición requiere un comentario no vacío.

\section{Límites declarados}

\begin{itemize}
  \item El WebSocket depende de una configuración compatible de Channels y Redis cuando hay más de
  un proceso.
  \item La entrega real de correo depende del proveedor y de un remitente validado.
  \item Supabase Storage funciona únicamente cuando sus credenciales y bucket están configurados.
  \item Existe una línea base de cobertura crítica, pero todavía no un umbral bloqueante ni un E2E
        completo en navegador.
  \item La licencia, cesión y soporte deben formalizarse por escrito antes de una transferencia final.
\end{itemize}

\ManualNote{La fuente canónica de capacidades es el código y sus pruebas. Este manual resume; no
reemplaza los permisos del backend ni los procedimientos de la plataforma de despliegue.}
